\section{Forecasting Performance}
\label{sec:forecasting_performance}

This section evaluates the forecasting behaviour of the proposed dCeNN--ELM backbone in comparison with the two selected baselines, namely Ridge Regression and LSTM. Since the anomaly-detection pipeline developed in this thesis is residual-based, forecasting quality is not an isolated objective; it directly determines the statistical reliability of downstream anomaly generation. For this reason, forecasting performance is assessed at two levels: first, through a benchmark-level comparison across models; and second, through a more detailed examination of target-wise and month-wise prediction behaviour in the 2022 test period \cite{repo_thesis_grid_anomaly}.

A careful interpretation of the forecasting results is required. The purpose of this section is not merely to identify the numerically strongest regressor, but to understand how the proposed architecture behaves relative to the intended design goal of the thesis. In particular, the dCeNN--ELM model is expected to occupy a middle ground between the simplicity of Ridge Regression and the heavier sequential expressiveness of LSTM. The forecasting results therefore need to be interpreted jointly in terms of accuracy, computational efficiency, and suitability as a backbone for the later neuro-symbolic anomaly analysis.

\subsection{Benchmark-Level Comparison}

The benchmark comparison table available in the repository summarizes the main model-level results across training time, inference latency, and target-wise RMSE for solar, wind, load, and price \cite{repo_thesis_grid_anomaly}. These results are reproduced in Table~\ref{tab:benchmark_comparison_ch71}.

\begin{table}[H]
	\centering
	\caption{Benchmark comparison of the proposed dCeNN--ELM model against Ridge and LSTM on the committed repository results.}
	\label{tab:benchmark_comparison_ch71}
	\small
	\begin{tabular}{|p{2.5cm}|p{1.8cm}|p{2.0cm}|p{1.8cm}|p{1.8cm}|p{1.8cm}|p{1.8cm}|}
		\hline
		\textbf{Model} & \textbf{Train Time (s)} & \textbf{Inference Latency (ms)} & \textbf{RMSE Solar} & \textbf{RMSE Wind} & \textbf{RMSE Load} & \textbf{RMSE Price} \\
		\hline
		dCeNN--ELM (Proposed) & 1.642 & 0.0873 & 0.1024 & 0.3236 & 4663.77 & 219.98 \\
		\hline
		LSTM (Baseline) & 430.085 & 0.3284 & 0.0857 & 0.2263 & 6266.13 & 252.34 \\
		\hline
		Ridge (Linear) & 0.0418 & 0.0649 & 0.0851 & 0.0782 & 219.94 & 20.05 \\
		\hline
	\end{tabular}
\end{table}

The committed benchmark results reveal a mixed but highly informative pattern. Ridge Regression is the strongest model in terms of raw RMSE across all four targets and is also the fastest both in training and inference. LSTM, while much more expressive in principle, is substantially slower and does not outperform Ridge in the current benchmark snapshot. The proposed dCeNN--ELM model lies between these two baselines in computational behaviour: it is dramatically faster than LSTM in training and inference, but slower than Ridge, as expected for a nonlinear latent model. In forecasting accuracy, however, it does not surpass Ridge in the currently committed results and also remains weaker than LSTM for solar and wind RMSE.

This outcome is academically important and should be stated honestly. The proposed architecture does not dominate the baselines on pure forecasting error in the benchmark artifact currently stored in the repository. Nevertheless, the result is still meaningful for the thesis because the claimed value of the dCeNN--ELM model is not solely ``best RMSE at all costs.'' Rather, its intended role is to provide a lightweight nonlinear forecasting backbone that remains computationally modest, structurally interpretable, and easily integrable with threshold calibration, ASP refinement, and edge export. The benchmark comparison therefore establishes the trade-off context for the rest of the chapter: forecasting accuracy must be interpreted together with deployment and neuro-symbolic utility, not in isolation.

\FloatBarrier
\begin{figure}[H]
	\centering
	\includegraphics[width=\linewidth]{benchmark_comparison.png}
	\caption{Benchmark comparison across dCeNN--ELM, Ridge, and LSTM. This figure visualizes the committed repository comparison in terms of RMSE, training time, and inference latency.}
	\label{fig:benchmark_comparison_ch71}
\end{figure}
\FloatBarrier

\subsection{Detailed Forecasting Metrics on the 2022 Test Period}

The benchmark table provides a compact comparison across models, but the detailed forecasting behaviour of the anomaly pipeline is best understood through the dedicated test-period metrics reported in \texttt{metrics\_summary.csv}. These results summarize the forecasting quality of the model used in the anomaly-detection workflow across the four main targets: solar capacity factor, wind capacity factor, load, and price \cite{repo_thesis_grid_anomaly}. Table~\ref{tab:metrics_summary_ch71} reproduces the committed values.

\begin{table}[H]
	\centering
	\caption{Detailed forecasting metrics on the 2022 test set for the implemented anomaly-detection pipeline.}
	\label{tab:metrics_summary_ch71}
	\small
	\begin{tabular}{|p{2.2cm}|p{1.6cm}|p{1.6cm}|p{2.0cm}|p{1.5cm}|p{1.6cm}|p{2.5cm}|}
		\hline
		\textbf{Target} & \textbf{RMSE} & \textbf{MAE} & \textbf{MAPE (\%)} & \textbf{$R^2$} & \textbf{Points} & \textbf{Anomaly Count} \\
		\hline
		CF\_Solar & 0.1226 & 0.1085 & 1599.10 & -2.4313 & 8688 & 3502 \\
		\hline
		CF\_Wind & 0.2554 & 0.2214 & 1485.64 & -0.3206 & 8688 & 4155 \\
		\hline
		Load\_MW & 1934.12 & 1605.59 & 22.37 & -1.3213 & 8688 & 1237 \\
		\hline
		Price & 209.34 & 174.92 & 62.64 & -2.0673 & 8688 & 4927 \\
		\hline
	\end{tabular}
\end{table}

Several observations follow from these results.

First, forecasting difficulty is clearly target-dependent. The renewable capacity-factor targets exhibit relatively small RMSE values in absolute terms because they are normalized variables, whereas load and price show substantially larger absolute errors due to their scale and volatility. This is expected in principle, but the anomaly counts indicate that price and wind are especially challenging from the perspective of residual-based abnormality.

Second, the reported negative $R^2$ values across the four targets indicate that, in the current test-period snapshot, the implemented forecasting backbone does not outperform a naive mean predictor on these evaluation metrics. This should be acknowledged explicitly. Accordingly, the present results do not support a claim of standalone predictive superiority for dCeNN--ELM over strong tabular baselines such as Ridge Regression. However, the value of the proposed model should be interpreted more narrowly and more appropriately. In the committed benchmark results, dCeNN--ELM remains substantially lighter than the LSTM baseline in parameter count, training time, and inference latency, while also improving over LSTM on the operationally important load and price targets. Specifically, dCeNN--ELM uses 8{,}016 parameters, trains in 1.64\,s, and achieves an inference latency of 0.087\,ms, compared with 79{,}364 parameters, 430.08\,s, and 0.328\,ms for LSTM; in addition, it improves on LSTM in load RMSE (4663.77 vs.\ 6266.13) and price RMSE (219.98 vs.\ 252.34) \cite{repo_thesis_grid_anomaly}. Ridge Regression remains the strongest pure forecaster, which suggests that the current feature-engineering regime captures a large share of the predictive structure in a form that is already well exploited by a linear model. The contribution of dCeNN--ELM in this thesis therefore lies not in universal forecasting dominance, but in providing a compact learned nonlinear backbone that can be integrated with threshold calibration, symbolic veto reasoning, and downstream utility analysis within a full neuro-symbolic anomaly-analysis pipeline.

Third, the MAPE values are extremely large for the renewable capacity-factor targets. This is not surprising because MAPE can become unstable when the true values are close to zero, which occurs frequently in solar and wind generation, especially during low-output periods. For this reason, RMSE and MAE are much more informative than MAPE for interpreting renewable forecasting quality in this context.

\subsection{Temporal Variability Across the 2022 Test Year}

A major strength of the committed repository outputs is that forecasting quality is also reported month by month in \texttt{metrics\_monthly.csv} \cite{repo_thesis_grid_anomaly}. This allows the model to be evaluated not only in aggregate, but also in terms of seasonal variability. Such temporal decomposition is especially important in electricity systems, where weather, renewable generation, load dynamics, and price volatility are strongly season-dependent.

The monthly results show that forecasting difficulty is highly non-uniform across the year. Several patterns are especially notable.

First, \textbf{summer months are the most challenging period for price forecasting}. August 2022 exhibits the highest monthly price RMSE among all reported months, reaching approximately 357.05, with September and July also showing extremely high price error levels \cite{repo_thesis_grid_anomaly}. This suggests that the model struggles most in periods of strong market instability, which is plausible given the wider volatility conditions that affected electricity prices in 2022.

Second, \textbf{wind forecasting also deteriorates markedly in summer and early autumn}. Wind RMSE rises above 0.30 in July, August, and September, indicating substantial difficulty in modeling wind-related dynamics during these periods \cite{repo_thesis_grid_anomaly}. This is consistent with the later anomaly results, where wind is one of the signals with a high anomaly burden.

Third, \textbf{load forecasting is somewhat more stable than price forecasting but still exhibits substantial month-to-month variation}. The highest monthly load RMSE values appear in the colder and transition months such as March and December, while late summer months such as August and September are somewhat less severe in absolute load RMSE terms \cite{repo_thesis_grid_anomaly}. This suggests that the difficulty of load forecasting is driven by a mix of seasonal demand patterns and model mismatch rather than one single regime.

Fourth, \textbf{solar forecasting is affected strongly by low-output periods and seasonal context}. December produces an especially poor $R^2$ value for solar, while other months also show unstable relative error behaviour, again confirming that normalized renewable variables remain challenging when values cluster near zero.

These month-wise results are important because they show that a single global forecasting metric can conceal periods of major instability. For an anomaly-detection framework, this matters directly: anomaly density and threshold behaviour are likely to be affected by the seasonal structure of residuals. The monthly analysis therefore strengthens the case for later threshold calibration and symbolic refinement, especially in the most volatile seasons.

\FloatBarrier
\begin{figure}[H]
	\centering
	\includegraphics[width=\linewidth]{master_timeline_Load_MW.png}
	\caption{Master timeline for the Load\_MW target over the 2022 test period, illustrating the temporal evolution of actual and predicted load, residual behaviour, and the corresponding anomaly-flag structure used in the downstream neuro-symbolic analysis.}
	\label{fig:monthly_forecasting_performance_load_ch71}
\end{figure}
\FloatBarrier


\subsection{Interpretation Relative to the Proposed Thesis Contribution}

The forecasting results should now be interpreted in relation to the actual contribution claim of the thesis. If the thesis were claiming that dCeNN--ELM is the strongest standalone forecaster in absolute statistical terms, the committed benchmark results would not support that claim. Ridge is clearly stronger in the available benchmark table, and the detailed metrics show that forecasting performance remains difficult across several signals and seasons.

However, that is not the most defensible reading of the framework. The more accurate interpretation is that dCeNN--ELM is proposed as a \emph{lightweight nonlinear forecasting backbone} within a larger neuro-symbolic anomaly-analysis system. Under that interpretation, three observations remain strongly relevant.

First, the proposed model preserves a substantial \textbf{efficiency advantage over LSTM}. Its training time and inference latency are much smaller than the sequential deep baseline, which supports the deployment and edge-feasibility arguments developed earlier in the thesis \cite{repo_thesis_grid_anomaly}.

Second, the architecture remains \textbf{structurally compatible with modular refinement}. Its residuals, thresholds, anomaly flags, and later ASP reasoning stages are all cleanly integrated within the repository pipeline. This architectural role is more important than raw RMSE alone for the later neuro-symbolic contribution.

Third, the forecasting difficulty itself helps motivate the symbolic layer. The fact that raw residual behaviour is noisy and seasonally unstable strengthens the need for ASP-based plausibility filtering rather than weakening it. In other words, the current forecasting results make the later anomaly-refinement results more, not less, important.

\subsection{Synthesis}

The forecasting results show that Ridge Regression remains the strongest model in
terms of raw predictive accuracy, while the proposed dCeNN--ELM backbone retains
a substantial efficiency advantage over LSTM. The monthly analysis further shows
that forecasting difficulty varies considerably across the 2022 test period,
particularly for price and wind. These findings establish the forecasting and
residual context for the symbolic anomaly-refinement analysis presented in the
following section.


\section{Neuro-Symbolic Refinement and Alert Reduction}
\label{sec:ns_refinement_false_alarm_reduction}

This section evaluates how the ASP layer reshapes the raw anomaly candidates
produced by residual thresholding. The symbolic stage applies persistence,
physical-plausibility, calendar, and cross-signal conditions to determine which
candidates are retained or vetoed. The analysis therefore focuses on candidate
counts, retention rates, veto categories, and the temporal distribution of symbolic
refinement.


\subsection{Raw Statistical Anomaly Burden}

The raw anomaly burden can be read directly from the committed repository evaluation outputs. According to the global metrics summary for the 2022 test period, the raw statistical detector produced the following anomaly counts: 3502 for solar capacity factor, 4155 for wind capacity factor, 1237 for load, and 4927 for price \cite{repo_thesis_grid_anomaly}. These values indicate that the raw detector is especially active for price and wind, while load produces substantially fewer anomaly flags.

This pattern is analytically important. Price and wind are also among the most difficult signals in the forecasting evaluation, particularly during volatile months. Their large raw anomaly counts therefore suggest that forecasting instability and contextual uncertainty propagate directly into the raw anomaly set. This is precisely the scenario in which a symbolic plausibility layer becomes valuable: when the raw detector is informative but overly permissive.

To formalize the refinement process, let $N_j^{\mathrm{raw}}$ denote the number of raw anomaly flags for target $j$, and let $N_j^{\mathrm{ref}}$ denote the number of anomalies that remain after ASP refinement. Then the number of eliminated alarms is
\begin{equation}
	N_j^{\mathrm{veto}} = N_j^{\mathrm{raw}} - N_j^{\mathrm{ref}},
	\label{eq:veto_count_ch72}
\end{equation}
and the corresponding retention rate is
\begin{equation}
	R_j = \frac{N_j^{\mathrm{ref}}}{N_j^{\mathrm{raw}}}\times 100.
	\label{eq:retention_rate_ch72}
\end{equation}
These quantities provide the most direct summary of candidate reduction and
retention in the current neuro-symbolic design.

% Replace the placeholder values in the Refined/Vetoed/Retention columns
% using the final committed or generated outputs from anomalies_refined.csv
% and/or the veto analysis summary produced by the repository scripts.

\begin{table}[htbp]
	\centering
	\caption{Raw versus refined anomaly counts by target. Raw counts are taken from the repository metrics summary; refined counts, vetoed counts, and retention rates from the final refined anomaly outputs.}
	\label{tab:raw_vs_refined_counts_ch72}
	\small
	\begin{tabular}{|p{2.5cm}|r|r|r|r|}
		\hline
		\textbf{Target} & \textbf{Raw Count} & \textbf{Refined Count} & \textbf{Vetoed Count} & \textbf{Retention Rate (\%)} \\
		\hline
		Solar & 3502 & 476 & 3026 & 13.59 \\
		Wind  & 4155 & 2238 & 1917 & 53.86 \\
		Load  & 1237 & 330 & 907 & 26.68 \\
		Price & 4927 & 3028 & 1899 & 61.46 \\
		\hline
	\end{tabular}
\end{table}

Even before the refined counts are inserted, the raw count structure already reveals where symbolic filtering is likely to matter most. Since price and wind have the largest raw anomaly burden, they are the most plausible candidates for meaningful symbolic refinement. Load, by contrast, may benefit more selectively from calendar-aware veto logic than from broad statistical pruning.

\subsection{Veto Reasoning Analysis}

The repository includes a dedicated visualization script, \texttt{src/21\_plot\_veto\_analysis.py}, which generates a target-wise comparison between confirmed symbolic anomalies and vetoed anomalies \cite{repo_thesis_grid_anomaly}. This figure should be treated as the primary visual summary of neuro-symbolic refinement in the chapter because it directly illustrates how much of the raw anomaly burden is retained and how much is suppressed after ASP reasoning.

Conceptually, the veto reasoning plot is important for two reasons. First, it quantifies the effect of the symbolic layer in a way that aggregate forecasting metrics cannot. Second, it shifts the interpretation of anomaly detection away from a pure count-maximization mindset. A stronger detector is not one that simply flags more timestamps; rather, it is one that preserves the most plausible anomalies while eliminating those that are inconsistent with domain knowledge.

\FloatBarrier
\begin{figure}[H]
	\centering
	\includegraphics[width=0.75\linewidth]{veto_reasoning_analysis.png}
	\caption{Target-wise veto reasoning analysis illustrating the impact of the ASP-based neuro-symbolic refinement layer on the raw anomaly set. For each signal, the figure compares anomalies retained after symbolic validation with anomalies vetoed by rule-based plausibility checks, thereby highlighting the extent of candidate reduction achieved by the proposed
		framework.}
	\label{fig:veto_analysis_ch73}
\end{figure}
\FloatBarrier


\subsection{Veto Reason Summary}

A more detailed explanation of symbolic candidate reduction is provided through
the report-level veto summary generated by
\texttt{src/25\_report\_tables.py}
\cite{repo_thesis_grid_anomaly}. An important methodological nuance must be stated explicitly: the current repository does \emph{not} store explicit \texttt{veto\_reason} atoms in the committed refined anomaly CSV. Therefore, the report-table generator constructs a set of \emph{artifact-only proxy categories} based on target type, timestamp context, and supporting raw anomaly fields. These proxy categories are still useful for interpretation, but they should be described honestly as report-level explanatory summaries rather than direct solver-extracted labels.

The repository defines the following main veto-reason categories:
\begin{itemize}
	\item \textbf{PV-at-night}, used when solar anomalies occur in time windows consistent with absent radiation;
	\item \textbf{Wind-ramp-guard (proxy)}, used when wind-related anomalies are associated with unusually large short-term changes;
	\item \textbf{Holiday/weekend-exemption (proxy)}, used when load-type anomalies coincide with weekend or holiday-like calendar context;
	\item \textbf{Missing co-anomaly confirmation}, used especially for price anomalies lacking the supporting load anomaly relation;
	\item \textbf{Transient non-persistent spike}, used for isolated or weakly supported anomalies that fail persistence-style confirmation.
\end{itemize}

These categories are highly valuable for the thesis because they move the discussion
beyond a generic statement such as ``ASP removed some anomalies.'' Instead, they
provide a domain-grounded interpretation of \emph{why} anomalies were removed. In
other words, they connect false-alarm reduction to concrete symbolic logic.

% Replace the placeholder counts and shares with values from
% reports/tables/veto_reason_summary.csv after generation.

\FloatBarrier
\begin{table}[htbp]
	\centering
	\caption{Veto reason summary based on repository-defined proxy categories. Final counts and shares from the generated \texttt{veto\_reason\_summary.csv} output.}
	\label{tab:veto_reason_summary_ch72}
	\small
	\begin{tabular}{|p{4.2cm}|r|r|}
		\hline
		\textbf{Veto Reason} & \textbf{Veto Count} & \textbf{Share of Vetoed Candidates (\%)} \\
		\hline
		PV-at-night & 2206 & 28.47 \\
		Wind-ramp-guard (proxy) & 83 & 1.07 \\
		Holiday/weekend-exemption (proxy) & 186 & 2.40 \\
		Missing co-anomaly confirmation & 1869 & 24.12 \\
		Transient non-persistent spike & 3405 & 43.94 \\
		\hline
	\end{tabular}
\end{table}
\FloatBarrier


\subsection{ASP Impact Heatmap}

The repository also includes an optional additional visualization in \texttt{src/24\_additional\_visuals.py}, which produces \texttt{asp\_impact\_heatmap.png} \cite{repo_thesis_grid_anomaly}. This plot displays monthly confirmed and vetoed anomaly counts across the four main targets. While not strictly necessary, it can substantially strengthen the thesis chapter because it reveals whether symbolic refinement is temporally concentrated in specific parts of the year.

This temporal view is valuable for two reasons. First, it helps link refinement behaviour to the monthly forecasting difficulty already discussed in Section~7.1. Second, it can reveal whether the symbolic layer is particularly active during volatile seasons such as summer and early autumn. If so, that would support the argument that ASP refinement is especially useful when raw statistical detection is under the greatest stress.

\FloatBarrier
\begin{figure}[H]
	\centering
	\includegraphics[width=\linewidth]{asp_impact_heatmap.png}
	\caption{Monthly ASP impact heatmap comparing confirmed and vetoed anomalies for each target in 2022. The left panel visualizes the anomalies retained after symbolic validation, whereas the right panel shows anomalies rejected by the ASP-based refinement layer. This comparison highlights the temporal and target-specific role of symbolic
		reasoning in reducing and reshaping the raw candidate set.}
	\label{fig:veto_analysis_ch73}
\end{figure}
\FloatBarrier


\subsection{Interpretation of Symbolic Refinement}

Taken together, the results show that the raw detector produces a broad candidate
set, while the ASP layer applies explicit domain conditions to retain or veto
individual candidates. The resulting anomaly set is smaller and accompanied by
interpretable proxy categories describing the main forms of symbolic rejection.
In the absence of complete ground-truth labels, these results demonstrate
domain-informed alert reduction and improved traceability rather than proving that
every vetoed candidate is a false positive.

\subsection{Synthesis}

The ASP layer substantially changes the raw anomaly distribution by applying
persistence, plausibility, calendar, and cross-signal constraints. The following
section evaluates the downstream financial consequences of this more selective
anomaly stream.


\section{Financial Impact of Refined Anomaly Signals}
\label{sec:financial_impact_refined_signals}

This section evaluates the downstream consequences of symbolic refinement through
the finance-aware backtest described in Chapter~5. The raw neural and refined
neuro-symbolic anomaly streams are compared using cumulative utility, penalties,
and final net utility. The purpose is to determine whether the more selective
refined stream improves the utility--risk balance of the resulting signals
\cite{repo_thesis_grid_anomaly}.


\subsection{Utility Formulation}

Let $r_t$ denote the reward or gain associated with a useful anomaly signal at time $t$, and let $p_t$ denote the penalty associated with a false, weak, or otherwise undesirable signal. The net utility at time $t$ is then written as
\begin{equation}
	u_t = r_t - p_t,
	\label{eq:net_utility_ch73}
\end{equation}
and cumulative utility up to time horizon $T$ is
\begin{equation}
	U(T) = \sum_{t=1}^{T} u_t.
	\label{eq:cumulative_utility_ch73}
\end{equation}

To compare the raw neural signal stream with the refined neuro-symbolic stream, the cumulative utility gap can be expressed as
\begin{equation}
	\Delta U(T) = U_{\mathrm{NS}}(T) - U_{\mathrm{N}}(T),
	\label{eq:utility_gap_ch73}
\end{equation}
where $U_{\mathrm{NS}}(T)$ denotes cumulative utility after ASP refinement and $U_{\mathrm{N}}(T)$ denotes cumulative utility from the raw neural anomaly stream. A positive value of $\Delta U(T)$ indicates that symbolic refinement improves the downstream economic value of anomaly detection.

\subsection{Cumulative Utility Over Time}

The most direct way to assess financial usefulness is to inspect cumulative utility over the full evaluation horizon. The repository includes a dedicated plotting script, \texttt{src/23\_plot\_utility\_comparison.py}, which produces a temporal comparison of utility accumulation over time \cite{repo_thesis_grid_anomaly}. This figure is important because it shows whether the advantage of the refined signals is stable, temporary, or concentrated in specific periods of the year.



\FloatBarrier
\begin{figure}[H]
	\centering
	\includegraphics[width=\linewidth]{utility_waterfall_comparison.png}
	\caption{Cumulative utility over time for the evaluated anomaly-signal variants.}
	\label{fig:cumulative_utility_over_time_ch73}
\end{figure}
\FloatBarrier

The utility comparison indicates that the purely neural strategy outperforms the neuro-symbolic variant in terms of cumulative utility and final net profit. Consequently, the ASP layer cannot be interpreted here as a direct profit-improving refinement mechanism in aggregate financial terms. Nevertheless, the neuro-symbolic strategy exhibits a lower penalty burden, with the figure reporting a risk reduction of approximately EUR~1{,}432. This suggests that symbolic reasoning contributes primarily through \emph{penalty mitigation} and more conservative signal selection rather than through maximization of gross return. Thus, in the present results, the financial value of the ASP layer lies in improving the quality and risk profile of accepted anomaly signals, even though this refinement comes at the cost of reduced total utility.

\subsection{Neural versus Neuro-Symbolic Utility Comparison}

Figure~\ref{fig:neural_vs_neuro_symbolic_utility_ch73} provides an aggregate
comparison of final net utility and total penalties for the raw neural and refined
neuro-symbolic strategies.

\FloatBarrier
\begin{figure}[H]
	\centering
	\includegraphics[width=\linewidth]{utility_waterfall_comparison1.png}
	\caption{Aggregate utility comparison between the raw neural anomaly stream and the neuro-symbolically refined anomaly stream. The figure summarizes the final net profit and total penalty burden for both strategies, highlighting that the purely neural strategy achieves higher cumulative utility, while the neuro-symbolic strategy reduces penalty exposure and thus improves the risk profile of accepted anomaly signals.}
	\label{fig:neural_vs_neuro_symbolic_utility_ch73}
\end{figure}
\FloatBarrier



\subsection{Utility and Penalty Summary}

Table~\ref{tab:utility_penalty_summary_ch73} reports the gross utility, total
penalties, net utility, and relative difference for the raw neural and refined
neuro-symbolic strategies.

% Fill the final values from the generated finance backtest and report summary outputs.

\begin{table}[H]
	\centering
	\caption{Utility and penalty summary for raw neural and neuro-symbolically refined anomaly signals. The neuro-symbolic net utility is taken from the existing 2022 financial backtest.}
	\label{tab:utility_penalty_summary_ch73}
	\small % Sets a base font size before scaling
	\resizebox{\textwidth}{!}{%
		\begin{tabular}{|p{4.5cm}|r|r|r|r|}
			\hline
			\textbf{Signal Variant} & \textbf{Gross Util.} & \textbf{Total Pen.} & \textbf{Net Util.} & \textbf{Relative Improv. \%} \\
			\hline
			Raw neural anomaly signals & 79{,}900 & 3{,}800 & 76{,}100 & -- \\
			\hline
			Neuro-symbolically refined signals & 45{,}800 & 2{,}400 & 43{,}389.59 & -43.0 \\
			\hline
		\end{tabular}%
	}
\end{table}

Table~\ref{tab:utility_penalty_summary_ch73} shows that the purely neural strategy achieves the higher overall net utility, whereas the neuro-symbolically refined strategy produces a lower penalty burden. In particular, the symbolic refinement reduces financial risk by approximately EUR~1{,}432, but this improvement in selectivity is accompanied by a substantial reduction in aggregate utility. Thus, in the present result snapshot, the ASP layer improves the risk profile of the anomaly stream rather than maximizing total simulated return.

\subsection{Interpretation of the Financial Results}

The financial results reveal a trade-off between aggregate return and penalty
exposure. The raw neural strategy achieves the higher net utility, whereas the
neuro-symbolic strategy incurs a lower penalty burden. The ASP layer should
therefore be interpreted as a more conservative, risk-mitigating refinement
mechanism rather than as a direct profit-maximization mechanism in the present
experiment.

\subsection{Synthesis}

The finance-aware evaluation shows that symbolic refinement reduces penalty
exposure but also lowers aggregate simulated utility relative to the raw neural
strategy. The next section examines how the retained anomalies appear at the level
of individual high-severity events.

\section{Case Studies of High-Severity Events}
\label{sec:case_studies_high_severity_events}

Aggregate metrics are useful for assessing overall forecasting behaviour, anomaly burden, and utility trends, but they cannot fully reveal how the proposed neuro-symbolic framework behaves on concrete real-world events. For that reason, this section presents a set of focused case studies drawn from the event-level anomaly table generated by the repository pipeline. The purpose is to move from global summary to event-level interpretation and to show how forecasting error, anomaly persistence, ASP-based refinement, and downstream utility interact in specific high-severity situations.

The case studies are based on the event table \texttt{reports/tables/anomaly\_events\_2022.csv}, which is generated from the refined anomaly stream using the event-building script \texttt{src/10\_build\_event\_table.py} \cite{repo_thesis_grid_anomaly}. In this pipeline, contiguous anomaly intervals are grouped into events and described using structured event attributes such as \texttt{event\_id}, \texttt{signal\_id}, \texttt{start\_ts}, \texttt{end\_ts}, \texttt{duration\_hours}, \texttt{n\_points}, \texttt{peak\_score}, \texttt{mean\_score}, and \texttt{severity\_label}. This event abstraction is important because it allows anomalies to be interpreted as meaningful temporal episodes rather than as isolated point detections.

For the present section, four representative high-severity events are selected:
\begin{enumerate}
	\item one \textbf{price event},
	\item one \textbf{load event},
	\item one \textbf{renewable event},
	\item one \textbf{multi-signal event}.
\end{enumerate}
This design ensures that the case-study analysis spans market behaviour, demand-side behaviour, renewable-side behaviour, and cross-signal system stress.

\subsection{Selection Criteria}

The selected events were chosen from the repository event table using three practical criteria. First, the event had to be labeled \texttt{HIGH} severity. Second, it had to exhibit a sufficiently large \texttt{peak\_score} to justify focused interpretation. Third, preference was given to events with either longer duration or clearer narrative value, since such events are easier to interpret in a thesis context than very short single-point spikes.

Table~\ref{tab:selected_high_severity_events} summarizes the four selected events that will be examined in detail.

\FloatBarrier
\begin{table}[H]
	\centering
	\caption{Selected high-severity case-study events drawn from \texttt{anomaly\_events\_2022.csv}.}
	\label{tab:selected_high_severity_events}
	\scriptsize
	\setlength{\tabcolsep}{7pt}
	\renewcommand{\arraystretch}{1.15}
	\begin{tabular}{|c|p{1.35cm}|p{2.0cm}|p{2.0cm}|c|c|c|c|p{1.1cm}|}
		\hline
		\textbf{ID} & \textbf{Signal} & \centering\textbf{Start (UTC)} & \centering\textbf{End (UTC)} & \textbf{Dur.} & \textbf{Pts} & \textbf{Peak} & \textbf{Mean} & \centering\textbf{Severity} \tabularnewline
		\hline
		1197 & price     & 2022-08-23 22:00 & 2022-08-24 18:00 & 20.0 & 21 & 8.0685 & 3.6603 & HIGH \\
		\hline
		640  & load\_mw  & 2022-07-20 04:00 & 2022-07-20 06:00 & 2.0  & 3  & 4.2729 & 3.5036 & HIGH \\
		\hline
		420  & cf\_wind  & 2022-07-28 04:00 & 2022-07-28 19:00 & 15.0 & 16 & 4.5360 & 2.8874 & HIGH \\
		\hline
		881  & n\_targets& 2022-08-22 10:00 & 2022-08-22 14:00 & 4.0  & 5  & 2.1360 & 1.9859 & HIGH \\
		\hline
	\end{tabular}
\end{table}
\FloatBarrier

The table shows that the chosen events are not trivial outliers. The selected price and wind events are especially strong because they are both severe and persistent. The selected load event is shorter, but it is still highly suitable for analysis because its peak score is large and its temporal localization makes the anomaly behaviour easy to interpret. The multi-signal event is included because it represents simultaneous stress across targets and therefore provides the clearest event-level illustration of multivariate anomaly behaviour.

\subsection{Case Study A: High-Severity Price Event}

The first case study focuses on the most prominent price event selected from the repository outputs: event ID 1197, associated with the \texttt{price} signal. This event spans from 2022-08-23 22:00 to 2022-08-24 18:00, lasts 20 hours, contains 21 anomalous points, and reaches a peak score of 8.0685 with a mean score of 3.6603. Among the available price events, this is one of the strongest sustained episodes in the committed event table and is therefore highly suitable for illustrating the behaviour of the anomaly pipeline under prolonged market stress.

From a methodological perspective, this event is valuable because price is also the most difficult target in the forecasting analysis of Section~7.1. A high-severity price event therefore provides a natural stress test for the full framework. The key questions to examine in the corresponding timeline and comparison plots are:
\begin{itemize}
	\item whether the forecast diverges sharply from the realized price path;
	\item whether the raw neural anomaly stream becomes highly active during the episode;
	\item whether the ASP layer retains the event rather than vetoing it;
	\item whether the finance-aware utility trajectory changes meaningfully over the event window.
\end{itemize}

If the symbolic layer confirms rather than suppresses this event, that would support the claim that the framework is capable of preserving economically meaningful price disturbances rather than over-filtering them.

\FloatBarrier
\begin{figure}[H]
	\centering
	\includegraphics[width=\linewidth]{master_timeline_Price.png}
	\caption{Master timeline for the Price signal, including actual vs predicted values, residuals, raw vs symbolic anomaly flags, and cumulative utility.}
	\label{fig:case_price_master_timeline}
\end{figure}
\FloatBarrier

\subsection{Case Study B: High-Severity Load Event}

The second case study focuses on event ID 640, associated with the \texttt{load\_mw} signal. This event runs from 2022-07-20 04:00 to 2022-07-20 06:00, lasts 2 hours, includes 3 anomalous points, and reaches a peak score of 4.2729 with a mean score of 3.5036. Although shorter than the selected price and wind events, it is still a strong case because its severity is high and its short duration makes the temporal evolution easy to interpret.

Load events are especially interesting in this thesis because load is influenced not only by system conditions but also by calendar context, weekends, and holidays. Consequently, load anomalies are among the signals most likely to be modified by the symbolic layer when context indicates that a raw deviation may not be operationally meaningful. For this reason, the load case study provides a good test of whether the ASP component can distinguish between plausible and implausible load deviations.

In the final figure set, the most important aspects to highlight for this event are:
\begin{itemize}
	\item the magnitude and abruptness of the load residuals,
	\item the concentration of anomaly flags in a short time interval,
	\item whether the event is retained or partially suppressed after ASP refinement,
	\item whether the event coincides with a noticeable change in utility trajectory.
\end{itemize}

\FloatBarrier
\begin{figure}[H]
	\centering
	\includegraphics[width=\linewidth]{master_timeline_Load_MW.png}
	\caption{Master timeline for the Load\_MW signal, showing forecast deviation, anomaly flags, and cumulative utility context.}
	\label{fig:case_load_master_timeline}
\end{figure}
\FloatBarrier


\subsection{Case Study C: High-Severity Renewable Event}

The third case study examines event ID 420, associated with the \texttt{cf\_wind} signal. This event spans 2022-07-28 04:00 to 2022-07-28 19:00, lasts 15 hours, contains 16 anomalous points, and reaches a peak score of 4.5360 with a mean score of 2.8874. Among the renewable events visible in the committed event table, this is one of the strongest long-duration wind cases and is therefore an excellent representative renewable anomaly episode.

This event is particularly important because renewable anomalies are where the hybrid structure of the thesis is most conceptually meaningful. Forecasting renewable behaviour is difficult due to weather uncertainty, but symbolic reasoning can still apply plausibility constraints and persistence-based refinement. A strong wind case therefore allows the thesis to show how the system handles a high-severity renewable episode that is likely to be both meteorologically meaningful and operationally relevant.

The corresponding plots should be interpreted with attention to:
\begin{itemize}
	\item whether the forecast misses a sustained deviation in wind-related behaviour;
	\item whether the anomaly remains persistent over a long enough interval to justify event grouping;
	\item whether the symbolic layer confirms the event rather than treating it as a transient ramp artifact;
	\item whether the event contributes positively to the downstream utility analysis.
\end{itemize}

\FloatBarrier
\begin{figure}[H]
	\centering
	\includegraphics[width=0.9\linewidth]{master_timeline_CF_Wind.png}
	\caption{Master timeline for the CF\_Wind signal, showing actual vs predicted behaviour, anomaly flags, and utility context.}
	\label{fig:case_wind_master_timeline}
\end{figure}
\FloatBarrier


\subsection{Case Study D: Multi-Signal Event}

The fourth case study focuses on the explicitly multivariate event stream represented by \texttt{n\_targets}. The selected event is event ID 881, which spans 2022-08-22 10:00 to 2022-08-22 14:00, lasts 4 hours, contains 5 anomalous points, and reaches a peak score of 2.1360 with a mean score of 1.9859. This event is not tied to a single variable, but instead captures a period in which multiple signals are jointly anomalous.

This case is especially important for the thesis because it most directly reflects the multivariate character of the proposed anomaly framework. Many practically relevant system disturbances do not appear as isolated univariate deviations. Instead, they manifest as simultaneous irregularities across load, price, and renewable signals. The \texttt{n\_targets} event stream is therefore a natural way to represent cross-signal disturbance episodes.

A strong supporting visualization for this discussion is the joint anomaly intersection analysis produced by the additional visualization script \texttt{src/24\_additional\_visuals.py} \cite{repo_thesis_grid_anomaly}. That plot helps justify why a multi-signal case study is needed at all: it shows that the anomaly structure of the system includes meaningful overlaps among targets rather than only isolated per-signal events.

The multivariate case study should therefore emphasize:
\begin{itemize}
	\item simultaneous activation across targets,
	\item whether the symbolic layer preserves cross-signal consistency,
	\item whether the event aligns with strong market or system stress windows,
	\item whether the downstream utility analysis benefits from preserving this joint event.
\end{itemize}

\FloatBarrier
\begin{figure}[H]
	\centering
	\includegraphics[width=\linewidth]{joint_anomaly_upset.png}
	\caption{Supporting figure for multivariate anomaly overlap, motivating the multi-signal case study.}
	\label{fig:case_multisignal_support}
\end{figure}
\FloatBarrier


\subsection{Cross-Case Interpretation}

Taken together, the four selected events reveal several important properties of the full framework.

First, anomaly severity is highly heterogeneous across targets. The selected price and wind events are both persistent and extreme, indicating that some of the most important system disturbances are prolonged rather than isolated. By contrast, the selected load event shows that even short events can still be highly severe and worthy of detailed interpretation.

Second, the case studies reinforce the value of event-level abstraction. A simple point-anomaly view would treat each flagged timestamp independently, whereas the event table reveals duration, persistence, and internal severity structure. This makes the resulting anomaly analysis far more interpretable and much better suited to narrative discussion in a thesis.

Third, the multi-signal event confirms that the anomaly framework is not merely a collection of four separate univariate detectors. It is capable of surfacing periods of cross-target irregularity, which are arguably the most relevant situations from an operational point of view.

Finally, the case studies provide the strongest qualitative evidence for the overall thesis claim. The forecasting layer alone produces residuals; the anomaly layer converts those residuals into candidate signals; the symbolic layer refines them; and the finance-aware layer interprets their usefulness. In the case-study view, these layers can be seen working together on real episodes rather than only through aggregate metrics.

\subsection{Synthesis}

The selected cases illustrate sustained price disturbances, short high-severity
load deviations, prolonged renewable anomalies, and cross-signal events. Together,
they complement the aggregate results by showing how severity, persistence,
symbolic refinement, and utility context appear within concrete event windows.
These findings provide the basis for the broader interpretation presented in
Chapter~8.